\scanpage{6}
satisfaisant
\[
u_i(1)=1\qquad i=1,2,3
\]
et
\[
\boxed{xyz=1\quad\Longrightarrow\quad u_1(x)u_2(y)u_3(z)=1}\tag*{\unsure{(b)|(6)|($*$)}}
\]

Prouvons que les $u_i$ sont \deleted{des hom de} \inserted{égaux à un seul} et unique hom de groupes.

Faisons $x=1$, on trouve que $u_3(y^{-1})=u_2(y)^{-1}$

Donc $u_3(y)=u_2(y^{-1})^{-1}$, pour la $=$ raison $u_1(y)=u_2(y^{-1})^{-1}$, donc $u_1(y)=u_3(y)$ et de $=$ $u_1(y)=u_2(y)$. Donc $u_1=u_2=u_3$, et $u(y)=u(y^{-1})^{-1}$, donc $\deleted{\unsure{y|u}}u(y^{-1})=u(y)^{-1}$.

On tire maintenant de \unsure{a)|a,|($*$)}
\[
\begin{aligned}
u((xy)^{-1})^{-1}&=u(x)u(y)\\
\rotatebox{90}{$=$}&\\
u(xy)&
\end{aligned}
\]
donc \qquad $u(xy)=u(x)u(y)$ \qquad cqfd.

\uline{Question} Le foncteur $\varphi$ est il une équivalence ?
% archive sheet number: 5
